% Chunk: \S5.3, Causal Vector Fields.
% Source: Weinberg QFT Vol. I, physical PDF pp. 230--236 (printed pp. 207--213),
% ending immediately before the \S5.4 heading on p. 236.
% Equations: (5.3.1)--(5.3.46). Figures: none. Footnotes: none.
% Status: source-reviewed, compile-clean, and visually checked; frozen.
% Uncertainties: none.

\subsection{Causal Vector Fields}

We now take up the next simplest kind of field, which transforms as a
four-vector, the simplest non-trivial representation of the homogeneous
Lorentz group. There are massive particles, the $W^\pm$ and $Z^0$, that at
low energies are described by such fields and that play an increasing role in
modern elementary particle physics, so this example is not merely of
pedagogical interest. (Also, although we are here considering only massive
particles, one approach to quantum electrodynamics is to describe the photon
in terms of a massive vector field in the limit of very small mass). For the
moment we will suppose that only one species of particle is described by this
field (dropping the species label $n$); then we shall consider the possibility
that the field describes both a particle and a distinct antiparticle.

In the four-vector representation of the Lorentz group, the rows and columns
of the representation matrices $D(\Lambda)$ are labelled with four-component
indices $\mu,\nu$, etc., with
\begin{equation}
  D(\Lambda)^\mu{}_{\nu}=\Lambda^\mu{}_{\nu}.
  \tag{5.3.1}\label{eq:5.3.1}
\end{equation}
The annihilation and creation parts of the vector field are written:
\begin{equation}
  \phi^{(+)\mu}(x)=\sum_\sigma(2\pi)^{-3/2}\int \dd^3p\,
  u^\mu(\mathbf{p},\sigma)a_{\mathbf{p},\sigma}e^{\ii p\cdot x},
  \tag{5.3.2}\label{eq:5.3.2}
\end{equation}
\begin{equation}
  \phi^{(-)\mu}(x)=\sum_\sigma(2\pi)^{-3/2}\int \dd^3p\,
  v^\mu(\mathbf{p},\sigma)a^\dagger_{\mathbf{p},\sigma}e^{-\ii p\cdot x}.
  \tag{5.3.3}\label{eq:5.3.3}
\end{equation}
The coefficient functions $u^\mu(\mathbf{p},\sigma)$ and
$v^\mu(\mathbf{p},\sigma)$ for arbitrary momentum are given in terms of those
for zero momentum by Eqs.~\eqref{eq:5.1.21} and \eqref{eq:5.1.22}, which here read:
\begin{equation}
  u^\mu(\mathbf{p},\sigma)=(m/p^0)^{1/2}
  L(\mathbf{p})^\mu{}_{\nu}u^\nu(0,\sigma),
  \tag{5.3.4}\label{eq:5.3.4}
\end{equation}
\begin{equation}
  v^\mu(\mathbf{p},\sigma)=(m/p^0)^{1/2}
  L(\mathbf{p})^\mu{}_{\nu}v^\nu(0,\sigma).
  \tag{5.3.5}\label{eq:5.3.5}
\end{equation}
(We are using the usual summation convention for spacetime indices
$\mu,\nu$, etc.) Also, the coefficient functions at zero momentum are subject
to the conditions \eqref{eq:5.1.25} and \eqref{eq:5.1.26}:
\begin{equation}
  \sum_{\bar\sigma}u^\mu(0,\bar\sigma)J^{(j)}_{\bar\sigma\sigma}
  =\mathcal{J}^\mu{}_{\nu}u^\nu(0,\sigma)
  \tag{5.3.6}\label{eq:5.3.6}
\end{equation}
and
\begin{equation}
  -\sum_{\bar\sigma}v^\mu(0,\bar\sigma)J^{(j)*}_{\bar\sigma\sigma}
  =\mathcal{J}^\mu{}_{\nu}v^\nu(0,\sigma).
  \tag{5.3.7}\label{eq:5.3.7}
\end{equation}
The rotation generators $\boldsymbol{\mathcal{J}}^\mu{}_{\nu}$ in the
four-vector representation are given by Eq.~\eqref{eq:5.3.1} as
\begin{equation}
  (\mathcal{J}_k)^0{}_0=(\mathcal{J}_k)^0{}_i
  =(\mathcal{J}_k)^i{}_0=0,
  \tag{5.3.8}\label{eq:5.3.8}
\end{equation}
\begin{equation}
  (\mathcal{J}_k)^i{}_j=-\ii\epsilon_{ijk},
  \tag{5.3.9}\label{eq:5.3.9}
\end{equation}
with $i,j,k$ here running over the values 1, 2, and 3. We note in particular
that $\boldsymbol{\mathcal{J}}^2$ takes the form
\begin{equation}
  (\boldsymbol{\mathcal{J}}^2)^0{}_0
  =(\boldsymbol{\mathcal{J}}^2)^0{}_i
  =(\boldsymbol{\mathcal{J}}^2)^i{}_0=0,
  \tag{5.3.10}\label{eq:5.3.10}
\end{equation}
\begin{equation}
  (\boldsymbol{\mathcal{J}}^2)^i{}_j=2\delta^i{}_j.
  \tag{5.3.11}\label{eq:5.3.11}
\end{equation}
From Eqs.~\eqref{eq:5.3.6} and \eqref{eq:5.3.7} it follows then that
\begin{equation}
  \sum_{\bar\sigma}u^0(0,\bar\sigma)
  (\mathbf{J}^{(j)})^2_{\bar\sigma\sigma}=0,
  \tag{5.3.12}\label{eq:5.3.12}
\end{equation}
\begin{equation}
  \sum_{\bar\sigma}u^i(0,\bar\sigma)
  (\mathbf{J}^{(j)})^2_{\bar\sigma\sigma}=2u^i(0,\sigma)
  \tag{5.3.13}\label{eq:5.3.13}
\end{equation}
and
\begin{equation}
  \sum_{\bar\sigma}v^0(0,\bar\sigma)
  (\mathbf{J}^{(j)*})^2_{\bar\sigma\sigma}=0,
  \tag{5.3.14}\label{eq:5.3.14}
\end{equation}
\begin{equation}
  \sum_{\bar\sigma}v^i(0,\bar\sigma)
  (\mathbf{J}^{(j)*})^2_{\bar\sigma\sigma}=2v^i(0,\sigma).
  \tag{5.3.15}\label{eq:5.3.15}
\end{equation}
Also, we recall the familiar result that
$(\mathbf{J}^{(j)})^2_{\bar\sigma\sigma}=j(j+1)\delta_{\bar\sigma\sigma}$.
From Eqs.~\eqref{eq:5.3.12}--\eqref{eq:5.3.15} we see that there are just two possibilities for
the spin of the particle described by the vector field: either $j=0$, for
which at $\mathbf{p}=0$ only $u^0$ and $v^0$ are non-zero, or else $j=1$
(so that $j(j+1)=2$), for which at $\mathbf{p}=0$ only the space-components
$u^i$ and $v^i$ are non-zero. Let us look in a little more detail at each of
these two possibilities.

\subsubsection*{Spin Zero}

By an appropriate choice of normalization of the fields, we can take the only
non-vanishing component of $u^\mu(0)$ and $v^\mu(0)$ to have the conventional
values:
\[
  u^0(0)=\ii(m/2)^{1/2},\qquad v^0(0)=-\ii(m/2)^{1/2}.
\]
(The label $\sigma$ here takes only the single value zero, and is therefore
dropped.) Then Eqs.~\eqref{eq:5.3.4} and \eqref{eq:5.3.5} yield for general momenta
\begin{equation}
  u^\mu(\mathbf{p})=\ii p^\mu(2p^0)^{-1/2}
  \tag{5.3.16}\label{eq:5.3.16}
\end{equation}
and
\begin{equation}
  v^\mu(\mathbf{p})=-\ii p^\mu(2p^0)^{-1/2}.
  \tag{5.3.17}\label{eq:5.3.17}
\end{equation}
The vector annihilation and creation fields here are nothing but the
derivatives of the scalar annihilation and creation fields $\phi^{(\pm)}$ for a
spinless particle that were defined in the previous section:
\begin{equation}
  \phi^{(+)\mu}(x)=\partial^\mu\phi^{(+)}(x),\qquad
  \phi^{(-)\mu}(x)=\partial^\mu\phi^{(-)}(x).
  \tag{5.3.18}\label{eq:5.3.18}
\end{equation}
It is obvious that the causal vector field for a spinless particle is also just
the derivative of the causal scalar field:
\begin{equation}
  \phi^\mu(x)=\phi^{(+)\mu}(x)+\phi^{(-)\mu}(x)
  =\partial^\mu\phi(x).
  \tag{5.3.19}\label{eq:5.3.19}
\end{equation}
Hence we need not explore this case any further here.

\subsubsection*{Spin One}

From Eqs.~\eqref{eq:5.3.6} and \eqref{eq:5.3.7} we see immediately that the vectors
$u^i(0,0)$ and $v^i(0,0)$ for $\sigma=0$ are in the 3-direction. By a suitable
normalization of the fields, we can take these vectors to have the values
\begin{equation}
  u^\mu(0,0)=v^\mu(0,0)=(2m)^{-1/2}
  \begin{bmatrix}0\\0\\0\\1\end{bmatrix}
  \tag{5.3.20}\label{eq:5.3.20}
\end{equation}
with four-vector components listed in the modern order $0,1,2,3$. To find the
other components, we use Eqs.~\eqref{eq:5.3.6}, \eqref{eq:5.3.7}, and \eqref{eq:5.3.9} to calculate the
effect of the raising and lowering operators $J^{(1)}_1\pm \ii J^{(1)}_2$ on
$u$ and $v$. This gives:
\begin{equation}
  u^\mu(0,+1)=-v^\mu(0,-1)
  =-\frac{1}{\sqrt2}(2m)^{-1/2}
  \begin{bmatrix}0\\1\\+\ii\\0\end{bmatrix},
  \tag{5.3.21}\label{eq:5.3.21}
\end{equation}
\begin{equation}
  u^\mu(0,-1)=-v^\mu(0,+1)
  =\frac{1}{\sqrt2}(2m)^{-1/2}
  \begin{bmatrix}0\\1\\-\ii\\0\end{bmatrix}.
  \tag{5.3.22}\label{eq:5.3.22}
\end{equation}
Applying Eqs.~\eqref{eq:5.3.4} and \eqref{eq:5.3.5} now yields
\begin{equation}
  u^\mu(\mathbf{p},\sigma)=v^{\mu*}(\mathbf{p},\sigma)
  =(2p^0)^{-1/2}e^\mu(\mathbf{p},\sigma)
  \tag{5.3.23}\label{eq:5.3.23}
\end{equation}
where
\begin{equation}
  e^\mu(\mathbf{p},\sigma)\equiv L^\mu{}_{\nu}(\mathbf{p})e^\nu(0,\sigma)
  \tag{5.3.24}\label{eq:5.3.24}
\end{equation}
with
\begin{equation}
  e^\mu(0,0)=\begin{bmatrix}0\\0\\0\\1\end{bmatrix},\quad
  e^\mu(0,+1)=-\frac{1}{\sqrt2}
  \begin{bmatrix}0\\1\\+\ii\\0\end{bmatrix},\quad
  e^\mu(0,-1)=\frac{1}{\sqrt2}
  \begin{bmatrix}0\\1\\-\ii\\0\end{bmatrix}.
  \tag{5.3.25}\label{eq:5.3.25}
\end{equation}
The annihilation and creation fields \eqref{eq:5.3.2} and \eqref{eq:5.3.3} here are
\begin{equation}
  \phi^{(+)\mu}(x)=\phi^{(-)\mu\dagger}(x)
  =(2\pi)^{-3/2}\sum_\sigma\int\frac{\dd^3p}{\sqrt{2p^0}}
  e^\mu(\mathbf{p},\sigma)a_{\mathbf{p},\sigma}e^{\ii p\cdot x}.
  \tag{5.3.26}\label{eq:5.3.26}
\end{equation}
The fields $\phi^{(+)\mu}(x)$ and $\phi^{(+)\nu}(y)$ of course commute (or
anticommute) for all $x$ and $y$, but $\phi^{(+)\mu}(x)$ and
$\phi^{(-)\nu}(y)$ do not. Their commutator (for bosons) or anticommutator
(for fermions) is
\begin{equation}
  [\phi^{(+)\mu}(x),\phi^{(-)\nu}(y)]_\mp
  =\int\frac{\dd^3p}{(2\pi)^3 2p^0}e^{\ii p\cdot(x-y)}\Pi^{\mu\nu}(\mathbf{p})
  \tag{5.3.27}\label{eq:5.3.27}
\end{equation}
where
\begin{equation}
  \Pi^{\mu\nu}(\mathbf{p})\equiv\sum_\sigma
  e^\mu(\mathbf{p},\sigma)e^{\nu*}(\mathbf{p},\sigma).
  \tag{5.3.28}\label{eq:5.3.28}
\end{equation}
A straightforward calculation using Eq.~\eqref{eq:5.3.25} shows that
$\Pi^{\mu\nu}(0)$ is the projection matrix on the space orthogonal to the
time-direction, and Eq.~\eqref{eq:5.3.24} then shows that $\Pi^{\mu\nu}(\mathbf{p})$
is the projection matrix on the space orthogonal to the four-vector $p^\mu$:
\begin{equation}
  \Pi^{\mu\nu}(\mathbf{p})=\eta^{\mu\nu}+p^\mu p^\nu/m^2.
  \tag{5.3.29}\label{eq:5.3.29}
\end{equation}
The commutator (or anticommutator) \eqref{eq:5.3.27} may then be written in terms of
the $\Delta_+$ function defined in the previous section, as
\begin{equation}
  [\phi^{(+)\mu}(x),\phi^{(-)\nu}(y)]_\mp
  =\left[\eta^{\mu\nu}-\frac{\partial^\mu\partial^\nu}{m^2}\right]
  \Delta_+(x-y).
  \tag{5.3.30}\label{eq:5.3.30}
\end{equation}
For our present purposes, the important thing about this expression is that
for $x-y$ space-like it does not vanish and is \textit{even} in $x-y$. We can
therefore repeat the reasoning of the previous section in seeking to construct
a causal field; we form a linear combination of annihilation and creation fields
\[
  v^\mu(x)=\kappa\phi^{(+)\mu}(x)+\lambda\phi^{(-)\mu}(x)
\]
for which, for $x-y$ spacelike,
\[
  [v^\mu(x),v^\nu(y)]_\mp
  =\kappa\lambda[1\mp1]
  \left[\eta^{\mu\nu}-\frac{\partial^\mu\partial^\nu}{m^2}\right]
  \Delta_+(x-y)
\]
and
\[
  [v^\mu(x),v^{\nu\dagger}(y)]_\mp
  =(|\kappa|^2\mp|\lambda|^2)
  \left[\eta^{\mu\nu}-\frac{\partial^\mu\partial^\nu}{m^2}\right]
  \Delta_+(x-y).
\]
In order for both to vanish for space-like $x-y$, it is necessary and
sufficient that the spin one particles be bosons and that
$|\kappa|=|\lambda|$. By a suitable choice of phase of the one-particle states
we can give $\kappa$ and $\lambda$ the same phase, so that $\kappa=\lambda$,
and then drop the common factor $\kappa$ by redefining the overall normalization
of the field. After all this, we find that the causal vector field for a
massive particle of spin one is
\begin{equation}
  v^\mu(x)=\phi^{(+)\mu}(x)+\phi^{(+)\mu\dagger}(x).
  \tag{5.3.31}\label{eq:5.3.31}
\end{equation}
We note that this is real:
\begin{equation}
  v^\mu(x)=v^{\mu\dagger}(x).
  \tag{5.3.32}\label{eq:5.3.32}
\end{equation}
However, if the particles it describes carry a non-zero value of some
conserved quantum number $Q$, then we cannot construct an interaction that
conserves $Q$ out of such a field. Instead, we must suppose that there is
another boson of the same mass and spin which carries an opposite value of
$Q$, and construct the causal field as
\begin{equation}
  v^\mu(x)=\phi^{(+)\mu}(x)+\phi^{(+)c\mu\dagger}(x)
  \tag{5.3.33}\label{eq:5.3.33}
\end{equation}
or in more detail
\begin{align}
  v^\mu(x)={}&(2\pi)^{-3/2}\sum_\sigma\int\frac{\dd^3p}{\sqrt{2p^0}}
  \left[e^\mu(\mathbf{p},\sigma)a_{\mathbf{p},\sigma}e^{\ii p\cdot x}
  \right.\notag\\[-2pt]
  &\left.\qquad\qquad\qquad
  +e^{\mu*}(\mathbf{p},\sigma)a^{c\dagger}_{\mathbf{p},\sigma}e^{-\ii p\cdot x}
  \right],
  \tag{5.3.34}\label{eq:5.3.34}
\end{align}
where the superscript $c$ indicates operators that create the antiparticle
that is charge-conjugate to the particle annihilated by $\phi^{(+)\mu}(x)$.
This again is a causal field, but no longer real. We can also use this formula
for the case of a purely neutral, spin one particle that is its own
antiparticle, by simply setting $a^c_{\mathbf{p}}=a_{\mathbf{p}}$. In either case,
the commutator of a vector field with its adjoint is
\begin{equation}
  [v^\mu(x),v^{\nu\dagger}(y)]
  =\left[\eta^{\mu\nu}-\frac{\partial^\mu\partial^\nu}{m^2}\right]
  \Delta(x-y),
  \tag{5.3.35}\label{eq:5.3.35}
\end{equation}
where $\Delta(x-y)$ is the function \eqref{eq:5.2.13}.

The real and complex fields we have constructed for a massive, spin one
particle satisfy interesting field equations. First, since $p^\mu$ in the
exponential in Eq.~\eqref{eq:5.3.26} satisfies $p^2=-m^2$, the field satisfies the
Klein--Gordon equation:
\begin{equation}
  (\Box-m^2)v^\mu(x)=0,
  \tag{5.3.36}\label{eq:5.3.36}
\end{equation}
just as for the scalar field. In addition, since Eq.~\eqref{eq:5.3.24} shows that
\begin{equation}
  e^\mu(\mathbf{p},\sigma)p_\mu=0,
  \tag{5.3.37}\label{eq:5.3.37}
\end{equation}
we now have another equation
\begin{equation}
  \partial_\mu v^\mu(x)=0.
  \tag{5.3.38}\label{eq:5.3.38}
\end{equation}
In the limit of small mass, Eqs.~\eqref{eq:5.3.36} and \eqref{eq:5.3.38} are just the equations
for the potential four-vector of electrodynamics in what is called Lorentz
gauge.

However, we cannot obtain electrodynamics from just any theory of massive
spin one particles by letting the mass go to zero. The trouble can be seen by
considering the rate of production of a spin one particle by an interaction
density $\mathcal{H}_I=J_\mu v^{\mu\dagger}$, where $J_\mu$ is an arbitrary
four-vector current. Squaring the matrix element and summing over the spin
$z$-components of the spin one particle gives a rate proportional to
\[
  \sum_\sigma|\langle J_\mu\rangle e^\mu(\mathbf{p},\sigma)^*|^2
  =\langle J_\mu\rangle\langle J_\nu\rangle^*\Pi^{\mu\nu}(\mathbf{p}),
\]
where $\mathbf{p}$ is the momentum of the emitted spin one particle, and
$\langle J_\mu\rangle$ is the matrix element of the current (say, at $x=0$)
between the initial and final states of all other particles. The term
$p^\mu p^\nu/m^2$ in $\Pi^{\mu\nu}(\mathbf{p})$ will, in general, cause the
emission rate to blow up when $m\to0$. The only way to avert this catastrophe
is to suppose that $\langle J_\mu\rangle p^\mu$ vanishes, which in coordinate
space is just the statement that the current $J^\mu$ must be \textit{conserved},
in the sense that $\partial_\mu J^\mu=0$. Indeed, the need for conservation of
the current can be seen by simply counting states. A massive spin one particle
has three spin states, which can be taken as the states with helicity $+1$, 0,
and $-1$, while any massless, spin one particle like the photon can only have
helicities $+1$ and $-1$: the current conservation condition just ensures that
the helicity zero states of the spin one particle are not emitted in the limit
of zero mass.

The inversions can be dealt with in much the same way as for the scalar field
discussed in the previous section. To evaluate the effect of space inversion,
we need a formula for $e^\mu(-\mathbf{p},\sigma)$. Using
$L^\mu{}_{\nu}(-\mathbf{p})=\mathcal{P}^\mu{}_{\rho}
L^\rho{}_{\tau}(\mathbf{p})\mathcal{P}^\tau{}_{\nu}$ and Eq.~\eqref{eq:5.3.24}, we have
\begin{equation}
  e^\mu(-\mathbf{p},\sigma)=-\mathcal{P}^\mu{}_{\nu}
  e^\nu(\mathbf{p},\sigma).
  \tag{5.3.39}\label{eq:5.3.39}
\end{equation}
Also, to evaluate the effect of time-reversal we need a formula for
$(-1)^{1+\sigma}e^{\mu*}(-\mathbf{p},-\sigma)$. Using
$e^{\mu*}(-\sigma)=-e^\mu(\sigma)$ and the above formula for
$L^\mu{}_{\nu}(-\mathbf{p})$, we find
\begin{equation}
  (-1)^{1+\sigma}e^{\mu*}(-\mathbf{p},-\sigma)
  =\mathcal{P}^\mu{}_{\nu}e^\nu(\mathbf{p},\sigma).
  \tag{5.3.40}\label{eq:5.3.40}
\end{equation}
Using these results and the transformation properties of the annihilation and
creation operators given in Section 4.2, it is straightforward to work out the
inversion transformation properties of the annihilation and creation fields.
Once again we find that, in order for causal fields to be transformed into
other fields with which they commute at space-like separations, it is necessary
that the intrinsic space inversion, charge-conjugation, and time-reversal phases
for spin one particles and their antiparticles be related by
\begin{equation}
  \eta^c=\eta^*,
  \tag{5.3.41}\label{eq:5.3.41}
\end{equation}
\begin{equation}
  \xi^c=\xi^*,
  \tag{5.3.42}\label{eq:5.3.42}
\end{equation}
\begin{equation}
  \zeta^c=\zeta^*.
  \tag{5.3.43}\label{eq:5.3.43}
\end{equation}
(In particular all phases must be real if the spin one particle is its own
antiparticle.) With these phase conditions satisfied, our causal vector field
\eqref{eq:5.3.34} has the inversion transformation properties
\begin{equation}
  Pv^\mu(x)P^{-1}=-\eta^*\mathcal{P}^\mu{}_{\nu}
  v^\nu(\mathcal{P}x),
  \tag{5.3.44}\label{eq:5.3.44}
\end{equation}
\begin{equation}
  Cv^\mu(x)C^{-1}=\xi^*v^{\mu\dagger}(x),
  \tag{5.3.45}\label{eq:5.3.45}
\end{equation}
\begin{equation}
  Tv^\mu(x)T^{-1}=\zeta^*\mathcal{P}^\mu{}_{\nu}
  v^\nu(-\mathcal{P}x).
  \tag{5.3.46}\label{eq:5.3.46}
\end{equation}
In particular, the minus sign in Eq.~\eqref{eq:5.3.44} means that a vector field that
transforms as a polar vector, with no extra phases or signs accompanying the
matrix $\mathcal{P}^\mu{}_{\nu}$, describes a spin one particle with intrinsic
parity $\eta=-1$.
